\documentclass[12pt,a4paper]{article}

\usepackage[T2A]{fontenc}
\usepackage[utf8]{inputenc}
\usepackage[russian]{babel}
\usepackage{amsmath,amssymb,amsthm,mathtools}
\usepackage{array}
\usepackage{booktabs}
\usepackage{algorithm}
\usepackage{algpseudocode}
\usepackage{geometry}
\usepackage{enumitem}

\geometry{left=25mm,right=15mm,top=20mm,bottom=20mm}

\newtheorem{definition}{Определение}
\newtheorem{theorem}{Теорема}
\newtheorem{lemma}{Лемма}
\newtheorem{proposition}{Утверждение}
\newtheorem{example}{Пример}
\newtheorem{remark}{Замечание}

\DeclareMathOperator{\sgn}{sgn}
\DeclareMathOperator{\seq}{seq}
\DeclareMathOperator{\ord}{ord}

\title{Билет №10 \\ \small Представление о сетях Петри для анализа свойств поведения параллельных программ (разметка, функционирование, развертка, граф достижимости).  (СпБПУ)}
\author{Сериков Владислав Александрович}
\date{5 мая 2026}

\begin{document}

\maketitle
\tableofcontents
\newpage

\section{Сети Петри для анализа поведения параллельных программ}

\subsection{Мотивация: параллелизм, синхронизация и ресурсы}

При анализе параллельных программ необходимо описывать не только последовательность вычислений, но и следующие явления:
\begin{itemize}
    \item независимое выполнение нескольких процессов;
    \item синхронизацию процессов;
    \item конкуренцию за общий ресурс;
    \item взаимное исключение;
    \item блокировки и тупики;
    \item накопление заявок в очередях;
    \item достижимость нежелательных состояний.
\end{itemize}

Классические блок-схемы хорошо описывают последовательное управление, но неудобны для представления независимых и взаимодействующих параллельных процессов. Сети Петри дают графо-аналитический аппарат, в котором:
\begin{itemize}
    \item \emph{условия} моделируются позициями;
    \item \emph{события} моделируются переходами;
    \item наличие условия моделируется фишкой в позиции;
    \item изменение состояния системы моделируется срабатыванием перехода.
\end{itemize}

Сети Петри особенно удобны для описания асинхронного параллелизма: независимые события могут происходить в произвольном порядке, а синхронизация возникает автоматически в тех местах сети, где переход имеет несколько входных позиций.

\subsection{Неформальное описание сети Петри}

Сеть Петри представляет собой ориентированный двудольный граф двух типов вершин:
\begin{itemize}
    \item \textbf{позиции} --- изображаются кружками;
    \item \textbf{переходы} --- изображаются прямоугольниками или чертами.
\end{itemize}

Дуги соединяют только вершины разных типов:
\[
p \to t
\]
означает, что позиция $p$ является входной для перехода $t$, а дуга
\[
t \to p
\]
означает, что позиция $p$ является выходной для перехода $t$.

Позиции содержат фишки. Распределение фишек по позициям называется \textbf{разметкой} или \textbf{маркировкой}. Срабатывание переходов перемещает фишки по сети и тем самым изменяет состояние системы.

\subsection{Формальное определение сети Петри}

\begin{definition}[Сеть Петри]
Сетью Петри называется четверка
\[
C = (P,T,I,O),
\]
где:
\begin{itemize}
    \item $P = \{p_1,p_2,\dots,p_n\}$ --- конечное множество позиций;
    \item $T = \{t_1,t_2,\dots,t_m\}$ --- конечное множество переходов;
    \item $P \cap T = \varnothing$;
    \item $I : T \to 2^P$ --- входная функция, задающая множество входных позиций каждого перехода;
    \item $O : T \to 2^P$ --- выходная функция, задающая множество выходных позиций каждого перехода.
\end{itemize}
\end{definition}

Если $p \in I(t)$, то $p$ называется входной позицией перехода $t$.

Если $p \in O(t)$, то $p$ называется выходной позицией перехода $t$.

Для удобства также вводят обозначения:
\[
{}^\bullet t = I(t),
\qquad
t^\bullet = O(t),
\]
где ${}^\bullet t$ --- множество входных позиций перехода $t$, а $t^\bullet$ --- множество его выходных позиций.

Аналогично для позиции $p$:
\[
{}^\bullet p = \{t \in T \mid p \in O(t)\},
\qquad
p^\bullet = \{t \in T \mid p \in I(t)\}.
\]

Здесь ${}^\bullet p$ --- множество переходов, ведущих в позицию $p$, а $p^\bullet$ --- множество переходов, выходящих из позиции $p$.

\subsection{Взвешенные сети Петри}

В классической простой сети дуга либо существует, либо не существует. В более общем случае дуги могут иметь веса.

\begin{definition}[Взвешенная сеть Петри]
Взвешенной сетью Петри называется пятерка
\[
N = (P,T,F,W),
\]
где:
\begin{itemize}
    \item $P$ --- конечное множество позиций;
    \item $T$ --- конечное множество переходов;
    \item $F \subseteq (P \times T) \cup (T \times P)$ --- множество дуг;
    \item $W : F \to \mathbb{N}_{>0}$ --- функция весов дуг.
\end{itemize}
\end{definition}

Если дуга отсутствует, ее вес считается равным нулю. Для простых сетей все существующие дуги имеют вес $1$.

\subsection{Разметка сети Петри}

\begin{definition}[Разметка]
Разметкой сети Петри называется отображение
\[
M : P \to \mathbb{N}_0,
\]
где $M(p)$ --- количество фишек в позиции $p$.
\end{definition}

Если позиции упорядочены:
\[
P = \{p_1,p_2,\dots,p_n\},
\]
то разметку удобно записывать вектором:
\[
M = (M(p_1),M(p_2),\dots,M(p_n)).
\]

\begin{definition}[Маркированная сеть Петри]
Маркированной сетью Петри называется пара
\[
(N,M_0),
\]
где $N$ --- сеть Петри, а $M_0$ --- начальная разметка.
\end{definition}

\subsection{Разрешенность перехода}

\begin{definition}[Разрешенный переход]
Переход $t \in T$ называется разрешенным в разметке $M$, если во всех его входных позициях достаточно фишек:
\[
\forall p \in P:
\quad
M(p) \ge W(p,t).
\]
\end{definition}

В простой сети это условие принимает вид:
\[
\forall p \in {}^\bullet t:
\quad
M(p) \ge 1.
\]

То есть переход разрешен тогда и только тогда, когда в каждой его входной позиции есть хотя бы одна фишка.

\subsection{Срабатывание перехода}

Если переход $t$ разрешен в разметке $M$, он может сработать. Срабатывание перехода:
\begin{itemize}
    \item удаляет фишки из входных позиций;
    \item добавляет фишки в выходные позиции.
\end{itemize}

Формально новая разметка $M'$ определяется формулой:
\[
M'(p) = M(p) - W(p,t) + W(t,p),
\]
где:
\[
W(p,t)=0,
\quad \text{если дуги } p \to t \text{ нет,}
\]
и
\[
W(t,p)=0,
\quad \text{если дуги } t \to p \text{ нет.}
\]

Обозначение:
\[
M \xrightarrow{t} M'
\]
означает, что из разметки $M$ при срабатывании перехода $t$ получается разметка $M'$.

\subsection{Пример простой сети Петри}

Рассмотрим сеть:
\[
P = \{p_1,p_2,p_3\},
\qquad
T = \{t_1,t_2\}.
\]

Пусть:
\[
{}^\bullet t_1 = \{p_1\},
\qquad
t_1^\bullet = \{p_2\},
\]
\[
{}^\bullet t_2 = \{p_2\},
\qquad
t_2^\bullet = \{p_3\}.
\]

Начальная разметка:
\[
M_0 = (1,0,0).
\]

Тогда:
\[
t_1 \text{ разрешен в } M_0,
\]
поскольку в $p_1$ есть фишка.

После срабатывания $t_1$:
\[
M_0 = (1,0,0) \xrightarrow{t_1} M_1 = (0,1,0).
\]

Теперь разрешен переход $t_2$, и:
\[
M_1 = (0,1,0) \xrightarrow{t_2} M_2 = (0,0,1).
\]

Таким образом, сеть описывает последовательный процесс:
\[
p_1 \to p_2 \to p_3.
\]

\subsection{Матричное представление сети Петри}

Для алгоритмического анализа удобно использовать матричную форму.

Пусть:
\[
|P| = n,
\qquad
|T| = m.
\]

Определим входную матрицу:
\[
A^- = (a^-_{ij}),
\]
где
\[
a^-_{ij} = W(p_i,t_j).
\]

Определим выходную матрицу:
\[
A^+ = (a^+_{ij}),
\]
где
\[
a^+_{ij} = W(t_j,p_i).
\]

\begin{definition}[Матрица инцидентности]
Матрицей инцидентности сети Петри называется матрица
\[
A = A^+ - A^-.
\]
\end{definition}

Если переход $t_j$ срабатывает один раз, то изменение разметки равно $j$-му столбцу матрицы $A$.

Пусть $e_j$ --- единичный вектор длины $m$, соответствующий переходу $t_j$. Тогда:
\[
M' = M + A e_j.
\]

\subsection{Последовательности срабатываний}

\begin{definition}[Последовательность срабатываний]
Последовательность переходов
\[
\sigma = t_{i_1}t_{i_2}\dots t_{i_k}
\]
называется допустимой из разметки $M_0$, если существует цепочка разметок:
\[
M_0 \xrightarrow{t_{i_1}} M_1
\xrightarrow{t_{i_2}} M_2
\to \dots
\xrightarrow{t_{i_k}} M_k.
\]
\end{definition}

В этом случае пишут:
\[
M_0 \xrightarrow{\sigma} M_k.
\]

\begin{definition}[Вектор Париха]
Для последовательности срабатываний $\sigma$ вектором Париха называется вектор
\[
\Psi(\sigma) \in \mathbb{N}_0^m,
\]
где $j$-я компонента показывает, сколько раз в последовательности $\sigma$ сработал переход $t_j$.
\end{definition}

Например, если:
\[
\sigma = t_1t_2t_1t_3t_1,
\]
то для множества переходов $T=\{t_1,t_2,t_3\}$ имеем:
\[
\Psi(\sigma) = (3,1,1)^T.
\]

\subsection{Уравнение состояния}

\begin{theorem}[Уравнение состояния сети Петри]
Пусть в маркированной сети Петри из начальной разметки $M_0$ достижима разметка $M$ с помощью последовательности срабатываний $\sigma$:
\[
M_0 \xrightarrow{\sigma} M.
\]
Тогда выполняется уравнение:
\[
M = M_0 + A \Psi(\sigma),
\]
где $A$ --- матрица инцидентности сети, а $\Psi(\sigma)$ --- вектор Париха последовательности $\sigma$.
\end{theorem}

\begin{proof}
Докажем утверждение индукцией по длине последовательности $\sigma$.

\textbf{База индукции.}

Если длина последовательности равна нулю, то $\sigma = \varepsilon$, переходы не срабатывают и:
\[
M = M_0.
\]
При этом:
\[
\Psi(\varepsilon) = 0.
\]
Следовательно:
\[
M_0 + A\Psi(\varepsilon)
=
M_0 + A \cdot 0
=
M_0.
\]
Значит, утверждение верно для пустой последовательности.

\textbf{Индукционный переход.}

Пусть утверждение верно для всех последовательностей длины $k$.

Рассмотрим последовательность длины $k+1$:
\[
\sigma' = \sigma t_j,
\]
где $\sigma$ имеет длину $k$.

Пусть:
\[
M_0 \xrightarrow{\sigma} M_k
\]
и затем:
\[
M_k \xrightarrow{t_j} M_{k+1}.
\]

По предположению индукции:
\[
M_k = M_0 + A\Psi(\sigma).
\]

При срабатывании перехода $t_j$ разметка изменяется на $j$-й столбец матрицы $A$, то есть:
\[
M_{k+1} = M_k + A e_j.
\]

Подставим выражение для $M_k$:
\[
M_{k+1}
=
M_0 + A\Psi(\sigma) + A e_j.
\]

Вынесем $A$:
\[
M_{k+1}
=
M_0 + A(\Psi(\sigma)+e_j).
\]

Но:
\[
\Psi(\sigma t_j) = \Psi(\sigma)+e_j.
\]

Следовательно:
\[
M_{k+1}
=
M_0 + A\Psi(\sigma').
\]

Утверждение доказано.
\end{proof}

\begin{remark}
Уравнение состояния дает необходимое условие достижимости. Если разметка $M$ достижима из $M_0$, то обязательно существует вектор $x \in \mathbb{N}_0^m$ такой, что:
\[
M = M_0 + Ax.
\]
Однако обратное в общем случае неверно: наличие неотрицательного целочисленного решения этого уравнения еще не гарантирует, что существует допустимый порядок срабатывания переходов.
\end{remark}

\subsection{Достижимость}

\begin{definition}[Достижимая разметка]
Разметка $M$ называется достижимой из начальной разметки $M_0$, если существует допустимая последовательность переходов $\sigma$ такая, что:
\[
M_0 \xrightarrow{\sigma} M.
\]
\end{definition}

\begin{definition}[Множество достижимости]
Множеством достижимости маркированной сети Петри называется множество:
\[
R(N,M_0)
=
\{M \mid \exists \sigma: M_0 \xrightarrow{\sigma} M\}.
\]
\end{definition}

Задача достижимости состоит в ответе на вопрос:

\[
M \in R(N,M_0)?
\]

То есть требуется определить, можно ли из начального состояния системы попасть в заданное состояние.

Для параллельных программ это позволяет проверять:
\begin{itemize}
    \item может ли программа попасть в состояние ошибки;
    \item возможен ли тупик;
    \item может ли быть одновременно занято два несовместимых ресурса;
    \item достижим ли участок программы;
    \item может ли возникнуть переполнение буфера.
\end{itemize}

\subsection{Граф достижимости}

Если множество достижимости конечно, его можно представить в виде графа.

\begin{definition}[Граф достижимости]
Графом достижимости маркированной сети Петри называется ориентированный размеченный граф
\[
G_R = (V,E),
\]
где:
\begin{itemize}
    \item вершины $V$ соответствуют достижимым разметкам;
    \item дуга
    \[
    M \xrightarrow{t} M'
    \]
    существует тогда и только тогда, когда переход $t$ разрешен в $M$ и его срабатывание переводит сеть в $M'$.
\end{itemize}
\end{definition}

Граф достижимости является аналогом графа состояний конечного автомата. Отличие состоит в том, что сеть Петри естественно описывает параллелизм, а граф достижимости уже разворачивает все возможные варианты выполнения в виде пространства состояний.

\subsection{Алгоритм построения графа достижимости}

Пусть дана маркированная сеть Петри $(N,M_0)$.

\textbf{Алгоритм построения графа достижимости для ограниченной сети:}

\begin{enumerate}
    \item Создать начальную вершину, соответствующую разметке $M_0$.
    \item Поместить $M_0$ в очередь необработанных разметок.
    \item Пока очередь необработанных разметок не пуста:
    \begin{enumerate}
        \item Извлечь очередную разметку $M$.
        \item Найти все переходы, разрешенные в $M$.
        \item Для каждого разрешенного перехода $t$ вычислить новую разметку $M'$.
        \item Добавить дугу
        \[
        M \xrightarrow{t} M'.
        \]
        \item Если вершина $M'$ еще не была добавлена в граф, добавить ее и поместить $M'$ в очередь необработанных разметок.
    \end{enumerate}
    \item Когда очередь пуста, построение завершено.
\end{enumerate}

\subsection{Минимальный пример построения графа достижимости}

Рассмотрим сеть:
\[
P = \{p_1,p_2\},
\qquad
T = \{t_1,t_2\}.
\]

Пусть:
\[
{}^\bullet t_1 = \{p_1\},
\qquad
t_1^\bullet = \{p_2\},
\]
\[
{}^\bullet t_2 = \{p_2\},
\qquad
t_2^\bullet = \{p_1\}.
\]

Начальная разметка:
\[
M_0 = (1,0).
\]

В $M_0$ разрешен только переход $t_1$:
\[
(1,0) \xrightarrow{t_1} (0,1).
\]

В разметке $(0,1)$ разрешен только переход $t_2$:
\[
(0,1) \xrightarrow{t_2} (1,0).
\]

Следовательно, граф достижимости состоит из двух вершин:
\[
M_0=(1,0),
\qquad
M_1=(0,1),
\]
и двух дуг:
\[
M_0 \xrightarrow{t_1} M_1,
\qquad
M_1 \xrightarrow{t_2} M_0.
\]

Эта сеть ограничена, безопасна и не имеет тупиков.

\subsection{Дерево достижимости}

Если сеть неограниченна, граф достижимости может быть бесконечным. В этом случае используют дерево достижимости или дерево покрываемости.

\begin{definition}[Дерево достижимости]
Дерево достижимости --- это дерево, корнем которого является начальная разметка $M_0$, а потомки каждой вершины строятся по всем разрешенным переходам.
\end{definition}

Если в процессе построения возникает разметка, которая уже встречалась на пути или в дереве, дальнейшее раскрытие этой вершины можно не выполнять, поскольку дальнейшее поведение уже было рассмотрено.

Если же обнаруживается рост числа фишек в некоторой позиции по сравнению с ранее встреченной разметкой, то для обозначения потенциально неограниченного роста используется символ
\[
\omega.
\]

Символ $\omega$ означает ``сколь угодно большое число фишек''.

\subsection{Алгоритм построения дерева покрываемости}

\textbf{Алгоритм Карпа--Миллера в упрощенной форме:}

\begin{enumerate}
    \item Создать корневую вершину с начальной разметкой $M_0$.
    \item Пометить корень как необработанный.
    \item Пока есть необработанные вершины:
    \begin{enumerate}
        \item Выбрать необработанную вершину с разметкой $M$.
        \item Если на пути от корня до $M$ уже встречалась такая же разметка, объявить вершину дублирующей и не раскрывать ее.
        \item Если в $M$ нет разрешенных переходов, объявить вершину терминальной.
        \item Иначе для каждого разрешенного перехода $t$ вычислить разметку $M'$.
        \item Если на пути от корня до $M$ существует разметка $M_a$ такая, что
        \[
        M' \ge M_a
        \]
        покомпонентно и хотя бы в одной компоненте строго больше, то каждую выросшую компоненту заменить на $\omega$.
        \item Добавить полученную разметку как потомка вершины $M$.
    \end{enumerate}
\end{enumerate}

\subsection{Иллюстрация дерева покрываемости}

Рассмотрим сеть:
\[
P=\{p_1,p_2\},
\qquad
T=\{t\}.
\]

Пусть переход $t$ имеет одну входную позицию $p_1$ и две выходные позиции $p_1,p_2$:
\[
{}^\bullet t = \{p_1\},
\qquad
t^\bullet = \{p_1,p_2\}.
\]

Начальная разметка:
\[
M_0=(1,0).
\]

Переход $t$ разрешен, поскольку в $p_1$ есть фишка.

После первого срабатывания:
\[
(1,0) \xrightarrow{t} (1,1).
\]

После второго:
\[
(1,1) \xrightarrow{t} (1,2).
\]

После третьего:
\[
(1,2) \xrightarrow{t} (1,3).
\]

Видно, что число фишек в позиции $p_2$ неограниченно растет. Поэтому дерево покрываемости можно записать как:
\[
(1,0) \to (1,\omega).
\]

Символ $\omega$ показывает, что в позиции $p_2$ может накопиться произвольно много фишек.

\subsection{Развертка сети Петри}

Граф достижимости перечисляет состояния системы, но теряет явную информацию о причинной независимости событий: разные перестановки независимых переходов могут приводить к одной и той же разметке.

Для анализа параллельных программ часто полезна \textbf{развертка} сети Петри.

\begin{definition}[Развертка сети Петри]
Развертка сети Петри --- это ациклическая сеть, представляющая возможные частичные порядки событий исходной сети. В развертке:
\begin{itemize}
    \item каждое событие соответствует конкретному срабатыванию перехода исходной сети;
    \item каждая позиция соответствует конкретному экземпляру условия;
    \item причинные зависимости отображаются дугами;
    \item независимые события не упорядочиваются искусственно.
\end{itemize}
\end{definition}

Идея развертки состоит в том, чтобы не строить все линейные последовательности срабатываний, а представить параллельное поведение как частично упорядоченное множество событий.

\subsection{Причинность, конфликт и параллелизм в развертке}

В развертке различают три отношения между событиями.

\begin{definition}[Причинность]
Событие $e_1$ причинно предшествует событию $e_2$, если результат $e_1$ является необходимым условием для $e_2$.
\end{definition}

Обозначение:
\[
e_1 < e_2.
\]

\begin{definition}[Конфликт]
События $e_1$ и $e_2$ находятся в конфликте, если они конкурируют за одну и ту же фишку, то есть срабатывание одного из них запрещает срабатывание другого.
\end{definition}

\begin{definition}[Параллелизм]
События $e_1$ и $e_2$ параллельны, если они не связаны причинностью и не находятся в конфликте.
\end{definition}

Развертки особенно полезны для уменьшения комбинаторного взрыва состояний. Например, если два перехода независимы, то в графе достижимости будут рассмотрены две последовательности:
\[
t_1t_2
\quad \text{и} \quad
t_2t_1,
\]
хотя фактически они описывают один и тот же параллельный сценарий. Развертка представляет этот случай как два независимых события.

\subsection{Пример параллельных независимых процессов}

Пусть имеются две независимые ветви:
\[
p_1 \xrightarrow{t_1} p_2,
\qquad
p_3 \xrightarrow{t_2} p_4.
\]

Начальная разметка:
\[
M_0=(1,0,1,0).
\]

В $M_0$ разрешены оба перехода:
\[
t_1, t_2.
\]

Возможны две последовательности:
\[
t_1t_2
\quad \text{и} \quad
t_2t_1.
\]

Обе приводят к одной разметке:
\[
(0,1,0,1).
\]

В графе достижимости эти варианты обычно представлены как разные пути, а в развертке события $t_1$ и $t_2$ будут показаны как независимые.

\subsection{Моделирование синхронизации}

Синхронизация двух процессов моделируется переходом с несколькими входными позициями.

Пусть:
\[
{}^\bullet t = \{p_1,p_2\}.
\]

Тогда переход $t$ разрешен только при условии:
\[
M(p_1)\ge 1
\quad \text{и} \quad
M(p_2)\ge 1.
\]

Это означает, что событие $t$ произойдет только тогда, когда оба процесса достигнут соответствующих точек синхронизации.

\subsection{Пример синхронизации}

Пусть первый процесс после некоторого вычисления помещает фишку в $p_1$, а второй процесс --- в $p_2$. Общий переход $t$ имеет входы $p_1$ и $p_2$ и выход $p_3$.

Если:
\[
M=(1,0,0),
\]
то $t$ не разрешен.

Если:
\[
M=(0,1,0),
\]
то $t$ также не разрешен.

Если:
\[
M=(1,1,0),
\]
то $t$ разрешен, и:
\[
(1,1,0) \xrightarrow{t} (0,0,1).
\]

Таким образом, сеть реализует операцию ожидания обоих параллельных процессов.

\subsection{Моделирование взаимного исключения}

Один из наиболее важных случаев анализа параллельных программ --- взаимное исключение при доступе к общему ресурсу.

Пусть два процесса имеют критические секции. Для ресурса вводится специальная позиция $r$:
\[
M(r)=1
\]
означает, что ресурс свободен.

Чтобы процесс вошел в критическую секцию, его переход входа должен потребить фишку из позиции $r$. При выходе из критической секции фишка возвращается в $r$.

\subsubsection{Схема}

Для первого процесса:
\[
p_1 + r \xrightarrow{t_1} c_1,
\]
\[
c_1 \xrightarrow{t_2} p_2 + r.
\]

Для второго процесса:
\[
q_1 + r \xrightarrow{u_1} c_2,
\]
\[
c_2 \xrightarrow{u_2} q_2 + r.
\]

Здесь:
\begin{itemize}
    \item $c_1$ --- критическая секция первого процесса;
    \item $c_2$ --- критическая секция второго процесса;
    \item $r$ --- свободный ресурс.
\end{itemize}

Начальная разметка:
\[
M_0(p_1)=1,
\qquad
M_0(q_1)=1,
\qquad
M_0(r)=1.
\]

Если первым сработает $t_1$, то фишка из $r$ исчезнет, и второй процесс не сможет войти в критическую секцию. После выхода первого процесса фишка возвращается в $r$, и ресурс снова становится доступным.

\subsection{Свойства сетей Петри}

\subsubsection{Безопасность}

\begin{definition}[Безопасная позиция]
Позиция $p$ называется безопасной, если во всех достижимых разметках:
\[
M(p) \le 1.
\]
\end{definition}

\begin{definition}[Безопасная сеть]
Сеть Петри называется безопасной, если каждая ее позиция безопасна.
\end{definition}

Безопасная позиция соответствует булевскому условию:
\[
M(p)=0 \quad \text{--- условие ложно,}
\]
\[
M(p)=1 \quad \text{--- условие истинно.}
\]

\subsubsection{Ограниченность}

\begin{definition}[$k$-ограниченная позиция]
Позиция $p$ называется $k$-ограниченной, если во всех достижимых разметках:
\[
M(p) \le k.
\]
\end{definition}

\begin{definition}[Ограниченная сеть]
Сеть Петри называется ограниченной, если существует число $k \in \mathbb{N}$ такое, что каждая позиция сети $k$-ограничена.
\end{definition}

Безопасность является частным случаем ограниченности при $k=1$.

\subsubsection{Сохраняемость}

\begin{definition}[Сохраняемая сеть]
Сеть Петри называется сохраняемой, если во всех достижимых разметках сохраняется общее число фишек:
\[
\sum_{p\in P} M(p) = \sum_{p\in P} M_0(p).
\]
\end{definition}

Иногда используется более общий вариант --- взвешенная сохраняемость.

\begin{definition}[Взвешенная сохраняемость]
Сеть Петри называется взвешенно сохраняемой, если существует вектор
\[
y \in \mathbb{N}_{>0}^{|P|}
\]
такой, что для любой достижимой разметки $M$:
\[
y^T M = y^T M_0.
\]
\end{definition}

\subsubsection{Живость}

\begin{definition}[Живой переход]
Переход $t$ называется живым, если из любой достижимой разметки можно попасть в такую разметку, где $t$ будет разрешен.
\end{definition}

\begin{definition}[Живая сеть]
Сеть Петри называется живой, если каждый ее переход живой.
\end{definition}

Живость означает отсутствие окончательного ``умирания'' частей системы. Для параллельных программ это связано с отсутствием ситуаций, когда некоторый фрагмент программы навсегда становится недоступным.

\subsubsection{Тупик}

\begin{definition}[Тупиковая разметка]
Разметка $M$ называется тупиковой, если в ней не разрешен ни один переход.
\end{definition}

То есть:
\[
\forall t\in T:
\quad t \text{ не разрешен в } M.
\]

Для параллельных программ тупиковая разметка часто соответствует deadlock-состоянию.

\subsection{Теорема о конечности графа достижимости}

\begin{theorem}
Пусть сеть Петри имеет конечное множество позиций $P$. Множество достижимых разметок конечно тогда и только тогда, когда сеть ограничена.
\end{theorem}

\begin{proof}
Обозначим:
\[
|P|=n.
\]

\textbf{Необходимость.}

Пусть множество достижимых разметок конечно:
\[
R(N,M_0)=\{M_1,M_2,\dots,M_s\}.
\]

Для каждой позиции $p_i$ рассмотрим множество значений:
\[
\{M_1(p_i),M_2(p_i),\dots,M_s(p_i)\}.
\]

Так как множество достижимых разметок конечно, это множество значений тоже конечно. Следовательно, для каждой позиции существует максимум:
\[
k_i = \max_{1\le j\le s} M_j(p_i).
\]

Положим:
\[
k = \max_{1\le i\le n} k_i.
\]

Тогда для любой достижимой разметки $M$ и любой позиции $p_i$:
\[
M(p_i) \le k.
\]

Следовательно, сеть ограничена.

\textbf{Достаточность.}

Пусть сеть ограничена. Тогда существует число $k$ такое, что для любой достижимой разметки $M$ и любой позиции $p_i$:
\[
M(p_i)\le k.
\]

Каждая разметка является вектором длины $n$:
\[
M=(M(p_1),M(p_2),\dots,M(p_n)).
\]

Каждая компонента этого вектора может принимать только одно из конечного числа значений:
\[
0,1,\dots,k.
\]

Следовательно, всего возможно не более:
\[
(k+1)^n
\]
различных разметок.

Значит, множество достижимых разметок конечно.
\end{proof}

\subsection{Инварианты сети Петри}

Инварианты позволяют доказывать свойства сети без полного построения графа достижимости.

\subsubsection{P-инварианты}

\begin{definition}[P-инвариант]
Вектор
\[
y \in \mathbb{Z}^{|P|}
\]
называется P-инвариантом, если:
\[
y^T A = 0.
\]
\end{definition}

\begin{theorem}[Сохранение P-инварианта]
Если $y$ --- P-инвариант сети Петри, то для любой достижимой разметки $M$:
\[
y^T M = y^T M_0.
\]
\end{theorem}

\begin{proof}
Пусть $M$ достижима из $M_0$ с помощью последовательности $\sigma$:
\[
M_0 \xrightarrow{\sigma} M.
\]

По уравнению состояния:
\[
M = M_0 + A\Psi(\sigma).
\]

Умножим обе части слева на $y^T$:
\[
y^T M = y^T M_0 + y^T A\Psi(\sigma).
\]

Так как $y$ является P-инвариантом:
\[
y^T A = 0.
\]

Следовательно:
\[
y^T A\Psi(\sigma)=0.
\]

Поэтому:
\[
y^T M = y^T M_0.
\]

Теорема доказана.
\end{proof}

\subsubsection{T-инварианты}

\begin{definition}[T-инвариант]
Вектор
\[
x \in \mathbb{Z}^{|T|}
\]
называется T-инвариантом, если:
\[
Ax = 0.
\]
\end{definition}

Если $x \ge 0$ и $Ax=0$, то срабатывание переходов в количествах, заданных вектором $x$, в сумме не меняет разметку:
\[
M' = M + Ax = M.
\]

T-инварианты описывают циклические режимы функционирования сети.

\subsection{Анализ взаимного исключения с помощью инварианта}

Рассмотрим модель двух процессов с общим ресурсом.

Пусть:
\[
r
\]
--- позиция свободного ресурса,
\[
c_1
\]
--- критическая секция первого процесса,
\[
c_2
\]
--- критическая секция второго процесса.

Если ресурс один, то начально:
\[
M_0(r)=1,
\qquad
M_0(c_1)=0,
\qquad
M_0(c_2)=0.
\]

Для корректной модели должен выполняться инвариант:
\[
M(r)+M(c_1)+M(c_2)=1.
\]

Это означает:
\begin{itemize}
    \item либо ресурс свободен;
    \item либо он занят первым процессом;
    \item либо он занят вторым процессом.
\end{itemize}

Из инварианта сразу следует невозможность одновременного нахождения обоих процессов в критических секциях:
\[
M(c_1)=1
\quad \text{и} \quad
M(c_2)=1
\]
невозможно, потому что тогда:
\[
M(r)+M(c_1)+M(c_2)\ge 2,
\]
что противоречит инварианту.

\subsection{Моделирование семафора}

Семафор можно представить позицией $s$.

Если семафор бинарный, то:
\[
M(s)=1
\]
означает, что семафор свободен,
а
\[
M(s)=0
\]
означает, что семафор занят.

Операция \texttt{wait} моделируется переходом, который потребляет фишку из $s$.

Операция \texttt{signal} моделируется переходом, который возвращает фишку в $s$.

Для счетного семафора с начальным значением $k$ задают:
\[
M_0(s)=k.
\]

Тогда одновременно ресурс могут захватить не более $k$ процессов.

\subsection{Моделирование производителя и потребителя}

Рассмотрим буфер ограниченной емкости $N$.

Введем позиции:
\begin{itemize}
    \item $empty$ --- количество свободных мест в буфере;
    \item $full$ --- количество занятых мест;
    \item $prod$ --- производитель готов произвести элемент;
    \item $cons$ --- потребитель готов потребить элемент.
\end{itemize}

Начальная разметка:
\[
M_0(empty)=N,
\qquad
M_0(full)=0.
\]

Переход производства:
\[
prod + empty \xrightarrow{produce} prod + full.
\]

Переход потребления:
\[
cons + full \xrightarrow{consume} cons + empty.
\]

Инвариант:
\[
M(empty)+M(full)=N.
\]

Он показывает, что сумма свободных и занятых ячеек буфера постоянна. Следовательно, переполнение буфера невозможно, если модель построена корректно.

\subsection{Пример анализа достижимости ошибки}

Пусть нужно проверить, может ли возникнуть ошибка взаимного исключения:
\[
M(c_1)=1,
\qquad
M(c_2)=1.
\]

Есть два подхода.

\subsubsection{Подход 1: граф достижимости}

\begin{enumerate}
    \item Построить граф достижимости из $M_0$.
    \item Проверить каждую достижимую разметку $M$.
    \item Если найдена разметка с:
    \[
    M(c_1)=1,
    \qquad
    M(c_2)=1,
    \]
    то ошибка достижима.
    \item Если такой разметки нет, взаимное исключение соблюдается.
\end{enumerate}

\subsubsection{Подход 2: инвариант}

Если доказан инвариант:
\[
M(r)+M(c_1)+M(c_2)=1,
\]
то состояние:
\[
M(c_1)=1,
\qquad
M(c_2)=1
\]
невозможно.

Этот способ часто эффективнее, потому что не требует полного перебора всех достижимых состояний.

\subsection{Основные задачи анализа сетей Петри}

К основным задачам анализа относятся:

\begin{enumerate}
    \item \textbf{Задача достижимости.}

    Даны $M_0$ и $M$. Требуется определить:
    \[
    M \in R(N,M_0)?
    \]

    \item \textbf{Задача ограниченности.}

    Требуется определить, существует ли $k$ такое, что:
    \[
    \forall M \in R(N,M_0),\ \forall p\in P:
    \quad
    M(p)\le k.
    \]

    \item \textbf{Задача безопасности.}

    Требуется определить:
    \[
    \forall M \in R(N,M_0),\ \forall p\in P:
    \quad
    M(p)\le 1.
    \]

    \item \textbf{Задача живости.}

    Требуется определить, может ли каждый переход потенциально срабатывать в дальнейшем из любой достижимой разметки.

    \item \textbf{Задача отсутствия тупиков.}

    Требуется проверить, существует ли достижимая разметка, в которой не разрешен ни один переход.

    \item \textbf{Задача сохраняемости.}

    Требуется проверить, сохраняется ли общее или взвешенное число фишек.
\end{enumerate}

\subsection{Классические сети Петри и время}

В классической сети Петри переходы считаются мгновенными. Это означает, что сеть описывает логическую структуру поведения, но не задает длительности операций.

Реальные действия часто имеют ненулевую длительность. Например:
\begin{itemize}
    \item выполнение команды процессором;
    \item ожидание ввода-вывода;
    \item передача сообщения;
    \item обслуживание заявки.
\end{itemize}

Такое длительное действие можно представить двумя мгновенными событиями:
\begin{itemize}
    \item начало действия;
    \item конец действия.
\end{itemize}

Между ними вводится позиция, означающая, что действие выполняется.

Например, выполнение операции процессором можно представить как:
\[
ready + cpu\_free \xrightarrow{start} running,
\]
\[
running \xrightarrow{finish} done + cpu\_free.
\]

Здесь позиция $running$ означает, что процессор занят выполнением операции.

\subsection{Пример модели простого компьютера}

Рассмотрим систему:
\begin{itemize}
    \item задания поступают во входную очередь;
    \item процессор обрабатывает одно задание;
    \item после обработки задание поступает на вывод.
\end{itemize}

Введем позиции:
\[
p_1 \text{ --- задание ожидает обработки,}
\]
\[
p_2 \text{ --- процессор свободен,}
\]
\[
p_3 \text{ --- задание выполняется,}
\]
\[
p_4 \text{ --- задание ожидает вывода.}
\]

Переходы:
\[
t_1 \text{ --- поступление задания,}
\]
\[
t_2 \text{ --- начало обработки,}
\]
\[
t_3 \text{ --- завершение обработки,}
\]
\[
t_4 \text{ --- вывод задания.}
\]

Связи:
\[
t_1^\bullet = \{p_1\},
\]
\[
{}^\bullet t_2 = \{p_1,p_2\},
\qquad
t_2^\bullet = \{p_3\},
\]
\[
{}^\bullet t_3 = \{p_3\},
\qquad
t_3^\bullet = \{p_4,p_2\},
\]
\[
{}^\bullet t_4 = \{p_4\}.
\]

Начальная разметка:
\[
M_0(p_2)=1,
\]
остальные позиции пусты.

Фишка в $p_2$ означает, что процессор свободен. Переход $t_2$ может сработать только тогда, когда есть задание в $p_1$ и процессор свободен. После начала обработки фишка из $p_2$ исчезает, поэтому другое задание не может одновременно занять процессор.

\subsection{Преимущества сетей Петри при анализе параллельных программ}

Сети Петри позволяют:
\begin{itemize}
    \item наглядно описывать параллельные процессы;
    \item моделировать синхронизацию;
    \item моделировать взаимное исключение;
    \item описывать очереди и ресурсы;
    \item формально определять множество достижимых состояний;
    \item обнаруживать тупики;
    \item доказывать безопасность с помощью инвариантов;
    \item строить графы достижимости и развертки;
    \item анализировать не только один сценарий выполнения, а все возможные сценарии.
\end{itemize}

\subsection{Ограничения классических сетей Петри}

Классические сети Петри имеют ряд ограничений:
\begin{itemize}
    \item переходы не имеют длительности;
    \item фишки не имеют индивидуальных данных;
    \item трудно напрямую моделировать сложные условия;
    \item для больших систем возникает комбинаторный взрыв числа состояний;
    \item бесконечные множества достижимости требуют специальных методов анализа.
\end{itemize}

Для преодоления этих ограничений используются расширения:
\begin{itemize}
    \item временные сети Петри;
    \item стохастические сети Петри;
    \item цветные сети Петри;
    \item иерархические сети Петри;
    \item E-сети;
    \item сети с ингибиторными дугами.
\end{itemize}

\subsection{Итоговая схема анализа параллельной программы с помощью сети Петри}

\begin{enumerate}
    \item Выделить процессы, ресурсы и события программы.
    \item Каждому условию или состоянию поставить в соответствие позицию.
    \item Каждому событию поставить в соответствие переход.
    \item Общие ресурсы представить специальными позициями с фишками.
    \item Синхронизацию представить переходами с несколькими входными позициями.
    \item Задать начальную разметку.
    \item Построить граф достижимости или дерево покрываемости.
    \item Проверить свойства:
    \begin{itemize}
        \item безопасность;
        \item ограниченность;
        \item отсутствие тупиков;
        \item живость;
        \item достижимость ошибочных состояний;
        \item корректность взаимного исключения.
    \end{itemize}
    \item При возможности найти P-инварианты и T-инварианты.
    \item Интерпретировать результаты анализа в терминах исходной программы.
\end{enumerate}

\subsection{Краткий вывод}

Сети Петри являются формальным и наглядным аппаратом для описания поведения параллельных программ. Основными понятиями являются позиции, переходы, фишки, разметки и правила срабатывания переходов. Динамика сети задается множеством достижимых разметок. Для конечных ограниченных сетей это множество представляется графом достижимости. Для неограниченных сетей применяются деревья покрываемости с символом $\omega$.

Сети Петри позволяют анализировать важные свойства параллельных программ: безопасность, ограниченность, живость, сохраняемость, достижимость и отсутствие тупиков. При этом граф достижимости дает полный перебор состояний для ограниченных сетей, а инварианты позволяют доказывать свойства системы более компактно.

\end{document}
